\documentclass[11pt]{article}
\usepackage[margin=1in]{geometry}
\usepackage{amsmath,amssymb}
\usepackage{booktabs}
\usepackage{hyperref}
\usepackage{array}
\usepackage[ruled,vlined]{algorithm2e}
\hypersetup{colorlinks=true,linkcolor=blue,citecolor=blue,urlcolor=blue}
\newcommand{\code}[1]{\texttt{#1}}
\title{A Nested Basket Model for the dunnhumby Complete Journey Data\\
\large Specification, inference, sampling, and an empirical study}
\author{}
\date{}
\begin{document}
\maketitle

\begin{abstract}
\noindent
We develop a model of supermarket baskets designed to answer price counterfactuals and
to generate synthetic shopping data. The model separates three decisions that a
single-choice model conflates --- whether a household buys from a category, which item
it then picks, and how much of it --- and estimates the coupling between them rather
than assuming it. Two commitments organise the work. First, every component is
introduced to answer a measurement taken on the data beforehand, and Section~\ref{sec:eda}
reports those measurements. Second, every claim is a number on held-out weeks, and where
a number is not comparable across settings we say so. The fitted model beats the
strongest simple baseline by $+0.342$ nats, its price coefficient retains 0.0\%
of its value when the price panel is scrambled before fitting, and baskets it generates
carry 87\% of the real price elasticity. Along the way we report a negative
result and its resolution: eleven changes to the model's interaction term failed to make
generated baskets reproduce which items co-occur, and a change to the \emph{negative
sampling distribution} succeeded.
\end{abstract}

\tableofcontents

%======================================================================
\section{Introduction}

\subsection{Main idea}

A retailer cutting one price wants to know what happens to the whole basket. Three
things can happen and they are economically distinct: households that were not going to
buy the category at all now do (\emph{incidence}); households already buying the
category switch to the discounted item (\emph{allocation}); and households buying that
item take more of it (\emph{quantity}). A model that only ranks items cannot tell these
apart, and the difference between them decides whether a promotion grows a category or
merely cannibalises it.

Our model gives each of these its own likelihood term, sharing one item utility
$u_{ijt}$. The coupling between incidence and allocation is a single estimated parameter
per category, $\kappa_c$, with a standard nested-logit reading: at $\kappa=1$ category
volume is whatever the items sum to, at $\kappa=0$ a price cut only moves share, and
above 1 the category expands more than proportionally.

A second idea runs through the paper. The purpose of the model is partly to
\emph{generate} data --- synthetic baskets on which a pricing or couponing policy could
be learned --- and generation is a far more demanding test than prediction. A model can
score well at ``which item completes this basket'' while producing synthetic baskets
that look nothing like real ones. We therefore evaluate the generator on five
model-free, distributional checks (Section~\ref{sec:gen}), and most of what we learned
came from those rather than from the likelihood.

\subsection{Main results}

\begin{enumerate}
\item \textbf{The three margins are separately identified and they are not equal.} A price
cut splits 79\% allocation, 12\% incidence, 9\% quantity (Section~\ref{sec:decomp}). The
quantity margin is invisible to a binary-purchase model.

\item \textbf{The price response survives a structural placebo.} Refitting on a scrambled
price panel retains 0.0\% of the coefficient while leaving $\kappa$ and the incidence
likelihood untouched (Section~\ref{sec:causal}). This rules out a spurious constant; it does
not identify the response against promotion timing, and we do not claim it does.

\item \textbf{It beats the published alternative on the published alternative's own
metric.} Against a reimplemented SHOPPER, $+0.320$ nats within category and $+0.310$ on the
full catalogue (Section~\ref{sec:strong}).

\item \textbf{The item conditional is exact.} Normalising over the purchased item's category
rather than the catalogue removes negative sampling entirely (Section~\ref{sec:negsamp}), so
no reported likelihood depends on a proposal distribution.

\item \textbf{Product embeddings carry real structure; household embeddings carry none.}
Sub-commodity purity at $27.8\times$ chance with two nulls at chance, against
gradient-boosted demographic probes that fail to beat the majority class on any target
(Section~\ref{sec:repr}).

\item \textbf{Matching means is not matching distributions.} A discriminator separates real
from generated baskets at AUC $0.77$, and the generator is 9\% short on basket size for
reasons we have not established (Section~\ref{sec:gen}).

\item \textbf{Two negative results and two reversals are reported rather than buried.}
Persistent contrastive divergence and an explicit habit feature both moved generation by
nothing; a temperature conclusion and a coupon-targeting conclusion both reversed once every
number was recomputed on a single fitted model (Sections~\ref{sec:gen}, \ref{sec:repr}).
\end{enumerate}

\subsection{Relation to prior work}

The model is a descendant of SHOPPER \cite{shopper} and of the nested factorization of
Donnelly et al.\ \cite{nf}. From SHOPPER we take the item-attribute factorisation, the
per-household price sensitivity as a low-rank bilinear form, and week effects as a
seasonality control. We depart from it in four ways: we add an explicit category-incidence stage
with an estimated nesting coefficient, we model units rather than binary purchase, we
condition on the store, and we replace its sampled softmax with an exact one over the
purchased item's category. Section~\ref{sec:related} discusses this in detail.

%======================================================================
\section{Data}
\label{sec:data}

The dunnhumby ``Complete Journey'' dataset: two years of transactions from one grocery
chain. Construction is in \code{scripts/22\_basket\_data.py}.

\begin{center}
\begin{tabular}{lr@{\qquad}lr}
\toprule
households $N$ & 2,066 & (basket, item) rows & 1,566,063 \\
items $J$ & 5,455 & baskets & 199,345 \\
categories $C$ & 188 & train / val / test rows & 1,228,695 / 134,665 / 202,703 \\
sub-commodities $S$ & 758 & units $>1$ & 22.3\% of rows, 42.6\% of units \\
stores & 115 & assortment density & 57.2\% \\
days $D$ & 712 & median repurchase gap & 28 days \\
\bottomrule
\end{tabular}
\end{center}

\paragraph{Filters.} An item needs at least 100 purchase lines, which is
statistical support for its own parameters; a household needs 20--300 trips.
Nothing is dropped to protect a modelling assumption --- in particular no basket is
discarded for containing two items from one category, which is the filter that
one-item-per-category models require and which removes most of the data.

\paragraph{Splits are temporal.} Train is week $<83$, validation 83--90,
test $\ge91$. A model whose purpose is to answer ``what happens next week'' must be
scored on later weeks; a random split would let it see each household's future.
Held-out rows whose item or household never appears in training are dropped rather than
scored as a cold start.

\paragraph{Prices.} The model sees a within-item log price deviation,
\begin{equation}
\Delta\log p_{jst} = \Delta\log p_{jt} + d_{jsw},
\label{eq:price}
\end{equation}
built in six steps set out in Section~\ref{sec:pricebuild}. Deviating within item means
the item's price \emph{level} is absorbed by its own intercept, so what identifies the
price response is an item moving against its own normal price rather than expensive items
differing from cheap ones.

\subsection{Notation}
\label{sec:notation}

\begin{center}
\begin{tabular}{ll@{\ }ll}
\toprule
symbol & meaning & symbol & meaning \\
\midrule
$i$ & household, $1..N$ & $N$ & 2,066 households \\
$j,k,m$ & item, $1..J$ & $J$ & 5,455 items \\
$c$ & category, $1..C$ & $C$ & 188 categories \\
$t$ & day, $1..D$ & $D$ & 712 days in the panel \\
$w$ & week-of-year, $1..52$ & $s$ & store, $1..115$ \\
$S$ & the set of items in a basket & $n$ & $|S|$, basket size \\
$B$ & the set of categories bought on a trip & $n_c$ & items of category $c$ the store carries \\
$R$ & purchase rows in a training batch & $C_r$ & candidate set for row $r$ \\
$\mathrm{sub}(j)$ & item $j$'s sub-commodity, of 758 & $\tau_{ijt}$ & days since $i$ last bought $\mathrm{sub}(j)$ \\
$\pi_j$ & within-category choice probability & $\lambda$ & a Poisson rate \\
\bottomrule
\end{tabular}
\end{center}

\subsection{How the price series is built}
\label{sec:pricebuild}

The price entering Eq.~\ref{eq:price} is the end of a chain, and each step is a choice.

\paragraph{Step 1: which price.} A dunnhumby transaction line records
four fields we need: \code{SALES\_VALUE} (what the retailer
banked), \code{QUANTITY}, \code{COUPON\_MATCH\_DISC} (the retailer's match of a
manufacturer coupon), and \code{RETAIL\_DISC} (the loyalty-card discount). Three different
prices can be reconstructed and they answer different questions. We use the
\textbf{loyalty price}, what a card holder faces at the shelf:
\begin{equation}
\mathrm{price}_{\ell} = \frac{\texttt{SALES\_VALUE} - \texttt{COUPON\_MATCH\_DISC}}{\texttt{QUANTITY}} .
\end{equation}
Subtracting the coupon match removes the part of the receipt the manufacturer refunded,
which is not a price the shopper responded to. We do \emph{not} use the paid price
(which adds \code{COUPON\_DISC} back), because two shoppers at the same shelf with
different coupons would then face different ``prices'' for the same item on the same day,
and the model has one price per (item, day).

\paragraph{Step 2: one number per item-day.} $\mathrm{price}_{jt}$ is the \emph{median}
loyalty price over all purchase lines of item $j$ on day $t$. A median rather than a mean
because a single mis-keyed quantity produces a wild per-unit value.

\paragraph{Step 3: fill the gaps.} Only 24.7\% of the $J\times D$ item-day grid is
directly observed. Missing days are filled by carrying the last observed price forward,
then the first observed price backward for days before an item's first sale. This assumes
a posted price persists until it is next seen to change, which is what a shelf tag does.

\paragraph{Step 4: log and deviate.}
\begin{equation}
\log p_{jt} = \log\max(\mathrm{price}_{jt},\,10^{-3}),
\qquad
\Delta\log p_{jt} = \log p_{jt} - \frac{1}{D}\sum_{t'=1}^{D}\log p_{jt'},
\end{equation}
where the mean runs over \textbf{all $D=712$ days of the panel}, filled days included,
so it is the item's average posted price over the two years rather than an average over
the days it happened to sell.

\paragraph{Step 5: the store deviation.} For each (item, store, week) actually observed,
\begin{equation}
d_{jsw} = \log \tilde p_{jsw} - \log \tilde p_{jw},
\end{equation}
where $\tilde p_{jsw}$ is the median loyalty price of item $j$ at store $s$ in week $w$
and $\tilde p_{jw}$ the median over the whole chain that week, kept only where
$|d_{jsw}|>10^{-6}$. There are 244,880 such cells, 0.53\% of the
(item-week $\times$ store) grid, so the overwhelming majority of lookups return
$d_{jsw}=0$ and fall back to the chain price. That is deliberate: imputing a store price
from nothing would invent variation the data does not contain.

\paragraph{Step 6: availability.} $\mathrm{carried}[j,s]$ is true if store $s$ sold item
$j$ \emph{at any point} in the two years. The threshold is one sale, not three: an item a
store ever sold was by definition available there, and a threshold of three marked
genuine purchases as unavailable, which set the inclusive value to $-10^9$ for categories
a household demonstrably bought from. The mask covers 57.2\% of the item $\times$ store
grid. dunnhumby has no stock-out feed, so this is a proxy and remains one.

%======================================================================
\section{What the data says}
\label{sec:eda}

Every measurement in this section was taken before any model was fitted, using counts
and least squares only. Each subsection ends with the design consequence, and
Section~\ref{sec:mapping} collects those consequences into a table. The code is
\code{scripts/29\_demand\_eda.py} and \code{scripts/30\_household\_eda.py}.

\subsection{Demand responds to price}
\label{sec:elast}

Form a panel with one row per (item, week). Let $b_{jw}$ be purchase lines and $T_w$ the
number of baskets chain-wide that week; put $y_{jw}=\log\frac{b_{jw}+0.5}{T_w}$, where
the $0.5$ keeps zero-purchase weeks in the panel, and let $x_{jw}$ be the mean log price.
The elasticity is a within-item least-squares slope, i.e.\ after removing each item's own
mean from both sides:
\begin{equation}
\hat\varepsilon=
\frac{\sum_{j,w}(x_{jw}-\bar x_j)(y_{jw}-\bar y_j)}{\sum_{j,w}(x_{jw}-\bar x_j)^2},
\qquad\text{556,410 item-weeks.}
\label{eq:within}
\end{equation}

\begin{center}
\begin{tabular}{lrl}
\toprule
outcome & elasticity & interpretation \\
\midrule
log buyers per trip & $-0.792$ & whether a household buys the item at all \\
log units per trip  & $-0.945$ & total units moved \\
log units per buyer & $-0.235$ & how much a buyer takes, given that they buy \\
\bottomrule
\end{tabular}
\end{center}

Of the total units response, 25\% comes from existing buyers taking more.
\textbf{Consequence:} a binary-purchase model discards a quarter of the price response,
so quantity needs its own likelihood term.

\subsection{A price cut moves three margins, and they can be measured separately}

Units from a category factor exactly,
\begin{equation}
\frac{\text{units}}{\text{basket}}
=\underbrace{\frac{\text{visits}}{\text{baskets}}}_{\text{incidence}}
\cdot\underbrace{\frac{\text{lines}}{\text{visit}}}_{\text{breadth}}
\cdot\underbrace{\frac{\text{units}}{\text{line}}}_{\text{depth}},
\end{equation}
so in logs the three elasticities must sum to the total. Estimating each by
Eq.~\ref{eq:within} on a (category, week) panel of 15,309 cells:

\begin{center}
\begin{tabular}{lrr}
\toprule
margin & elasticity & \\
\midrule
incidence & $-0.675$ & \\
breadth   & $-0.092$ & \\
depth     & $-0.158$ & \\
\midrule
sum & $-0.925$ & against $-0.945$ from Eq.~\ref{eq:within} on the item panel \\
\bottomrule
\end{tabular}
\end{center}

The two agree to 0.021 across different panels, different units of analysis and no
shared code. About 73\% of the entire units response is households deciding to
enter a category at all. \textbf{Consequence:} incidence is the dominant margin and must
be modelled explicitly, not folded into item choice.

\subsection{Promotions: an event study}

Take every item-week where log price fell by at least $0.15$ against the previous week
(37,132 events over 4,712 items) and normalise demand by the item's own mean.
Events are not independent: 57.7\% of consecutive cuts for the same item fall less
than seven weeks apart, so we also report the 20,725 \emph{isolated} events with no
other cut within $\pm3$ weeks. Neither $1.0$ nor the pre-period is a valid baseline ---
the former averages over the promotions themselves, the latter is selected to be
expensive by the event definition --- so we use a \emph{quiet week}, an item-week more
than three weeks from any cut of that item. There are 368,173, running at 0.896
of the item's mean demand at essentially its normal price.

\begin{center}
\begin{tabular}{lrrr}
\toprule
weeks from cut & demand & price (log dev.) & demand $\div$ quiet \\
\midrule
$-3$ & 0.925 & $+0.050$ & 1.03 \\
$-1$ & 0.870 & $+0.123$ & 0.97 \\
$\ \ 0$ & \textbf{2.136} & $-0.194$ & \textbf{2.38} \\
$+1$ & 1.337 & $-0.139$ & 1.49 \\
$+3$ & 1.065 & $-0.003$ & 1.19 \\
quiet & 0.896 & $+0.011$ & 1.00 \\
\bottomrule
\end{tabular}
\end{center}

Three readings. \emph{Demand does not ramp before a cut}: the pre-period sits within 3\%
of quiet, so the spike is not partly its own cause --- which is what licenses reading the
price coefficient causally, and is tested directly in Section~\ref{sec:causal}.
\emph{The promotion week is much steeper than elasticity}: from $-1$ to $0$, price moves
$-0.317$ and demand $+0.898$ in logs, an implied slope of $-2.83$, which is $3.4\times$
the $-0.792$ measured on ordinary variation --- a promoted price arrives with display and
feature that the price variable does not carry. \emph{Most of the tail is the discount
still running}: at $+1$ price is still $0.139$ below normal, and by $+3$, when price has
returned, demand is 1.19 of quiet. No offset dips below quiet, so no stockpiling is
visible within three weeks.

\subsection{Households differ, and by how much can be measured}
\label{sec:hh}

\paragraph{Taste.} Split each household's purchases in half by date, count purchases per
category in each half to get two vectors in $\mathbb{R}^{188}$, scale each to unit
length, and take the dot product. Against \emph{all} 4,266,290 cross-household pairs:

\begin{center}
\begin{tabular}{lr}
\toprule
own two halves & \textbf{0.7838} \\
two different households & 0.4710 \\
ratio & \textbf{1.66$\times$} \\
\bottomrule
\end{tabular}
\end{center}

A stranger already scores 0.4710 because everyone buys milk and bread; the gap is
what is personal. \textbf{Consequence:} a per-household taste vector $\theta_i$ is
justified.

\paragraph{Price sensitivity is real but poorly measured.} Regress $\log$ units on
$\log$ price per household, within (household, item), on pairs bought three or more
times. Across 1,640 households the median is $-0.169$ with p10--p90
$-0.427$--$-0.016$. But spread alone proves nothing --- noisy estimates are spread
out even when the truth is identical --- and measuring each household twice on disjoint
halves gives a split-half correlation of only $+0.236$ on 1,374 households.

The reason is visible in the estimator, which divides by the price variation that
household actually saw: 18.6\% of qualifying rows come from a (household, item) pair
whose price never moved. Splitting households by how much movement they had:

\begin{center}
\begin{tabular}{lrr}
\toprule
 & spread of slopes (sd) & share positive \\
\midrule
little price movement & 0.226 & 11.5\% \\
much price movement & 0.129 & 3.5\% \\
\bottomrule
\end{tabular}
\end{center}

Where prices moved, estimates are tighter and almost all negative. \textbf{Consequence:}
per-household price sensitivity is justified but must be low-rank rather than free, and a
targeting policy should weight by how much is known about each household.

\subsection{Stores are not interchangeable}

15.8\% of store-item-weeks sit more than a cent from the chain price, and the
median store ever sold only 63\% of the catalogue. \textbf{Consequence:} scoring
an item a store never stocked as a \emph{rejected} alternative is a specification error;
it was not available to reject.

\subsection{Base rates the model must reproduce}

Categories per basket 6.12 of 188 (3.25\%); items per basket 7.86;
units per basket 10.64; distinct items per purchased category 1.284. These are
the quantities Section~\ref{sec:gen} holds the generator to.

%======================================================================
\section{The model}
\label{sec:model}

A trip is one household $i$ shopping at store $s$ on day $t$, in week-of-year $w$. The
model answers four questions about that trip, in this order:

\begin{enumerate}
\item For each of the $C$ categories --- does the household buy from it at all?
\item For a category it buys from --- how many \emph{different} products does it take?
\item Which products?
\item How many units of each?
\end{enumerate}

Section~\ref{sec:util} builds one quantity that all four answers share.
Section~\ref{sec:heads} gives the four answers. Section~\ref{sec:nest} explains the one
term that ties question 1 to questions 2--4, and what its coefficient means.
Section~\ref{sec:inter} covers the term that lets items in a basket influence each other.
Section~\ref{sec:joint} multiplies the four answers into the probability of a whole basket
and says how the training objective differs from it.

\subsection{The shared quantity: how appealing item \texorpdfstring{$j$}{j} is}
\label{sec:util}

Everything rests on one number per (household, item, trip), which we call the item's
utility:
\begin{equation}
\begin{aligned}
u_{ijt}=\;&\underbrace{\lambda_j}_{\text{popular?}}
\;+\;\underbrace{\theta_i^{\top}\alpha_j}_{\text{their kind of thing?}}
\;+\;\underbrace{\rho_j^{\top}\bar\alpha_j(S)}_{\text{fits the cart?}}
\;-\;\underbrace{(\gamma_i^{\top}\beta_j)\,\Delta\log p_{jst}}_{\text{expensive today?}}\\[4pt]
&+\;\underbrace{\eta_j^{\top}x_{ijt}}_{\text{due a repurchase?}}
\;+\;\underbrace{\mu_j^{\top}\delta_w}_{\text{in season?}}
\;+\;\underbrace{\zeta_j^{\top}\xi_s}_{\text{this store's kind of thing?}}
\;+\;\underbrace{w^{\mathrm{loy}}_j\omega_{ij}+w^{\mathrm{frq}}_jf_{ij}}_{\text{their usual?}}
\end{aligned}
\label{eq:util}
\end{equation}

Each term is a dot product of two learned vectors, which is a compact way of writing a
large table of interactions with few numbers. $\theta_i^{\top}\alpha_j$ is an example:
rather than storing one taste value for every household--item pair ($N\times J$ numbers),
each household gets a vector $\theta_i$, each item a vector $\alpha_j$, and their dot
product supplies the pair. The same device gives every household its own price
sensitivity $\gamma_i^{\top}\beta_j$ instead of one shared slope.

\paragraph{The price input.} $\Delta\log p_{jst}$ is how far today's price sits from that
item's own normal price, in logs (Section~\ref{sec:pricebuild}). It enters with a minus
sign, so a \emph{positive} $\gamma_i^{\top}\beta_j$ means the household buys less when the
price is above normal.

\paragraph{The memory input.} $x_{ijt}$ describes how long it has been since this
household last bought something of this kind. Let $\tau$ be the days since household $i$
last bought item $j$'s sub-commodity; then
\begin{equation}
x_{ijt}=\bigl[\ \mathbb 1\{\text{never before}\},\quad e^{-\tau/7},\quad
e^{-\tau/g_{\mathrm{sub}(j)}},\quad \log(1+\tau)/\log 100\ \bigr]^{\top},
\end{equation}
with the last three set to zero if there is no earlier purchase. The three decays let the
model express ``bought it yesterday'', ``about due'', and ``long gone'' separately, and
$g_{\mathrm{sub}(j)}$ is that sub-commodity's median repurchase gap so that ``due'' means
a week for milk and months for shampoo.

\paragraph{The habit input.} $\theta_i^{\top}\alpha_j$ is a rank-$K$ bilinear form standing
in for an $N\times J$ table, so it can say ``this household likes this kind of product'' but
not ``this household buys \emph{this} product and no other''. Two scalars supply the
difference. With $n^{\mathcal T}_{ij}$ the number of training trips on which household $i$
bought item $j$,
\begin{equation}
\omega_{ij}=\frac{n^{\mathcal T}_{ij}}
{\max\bigl(1,\ \sum_{j'\in\mathcal{C}(c(j))}n^{\mathcal T}_{ij'}\bigr)},
\qquad
f_{ij}=\frac{\log(1+n^{\mathcal T}_{ij})}{\log 100}.
\label{eq:persist}
\end{equation}
$\omega_{ij}$ is the household's within-category share of item $j$ --- exactly the quantity
the item conditional normalises, so $\sum_{j\in\mathcal{C}(c)}\omega_{ij}=1$ for any
category the household has bought from --- and $f_{ij}$ the absolute strength of the habit.
Both are computed on the training window only, so they are past lookups rather than leakage,
and both are fixed \emph{data}: only the per-item loadings $w^{\mathrm{loy}}_j$,
$w^{\mathrm{frq}}_j$ are estimated. Being frozen at the training window, they do not update
inside a generated rollout, which is one reason Section~\ref{sec:sampling} produces a sampler
and not a simulator.

\subsection{The four answers}
\label{sec:heads}

Write $\mathcal{C}_s(c)$ for the items of category $c$ that store $s$ carries. All four
answers are read off $u_{ijt}$.

\paragraph{1. Is category $c$ bought?} A coin flip whose bias depends on the household,
the season of its last purchase, and how good the category looks today:
\begin{equation}
y_{ct}\sim\mathrm{Bernoulli}\bigl(\sigma(\eta_{ict})\bigr),\qquad
\eta_{ict}=\underbrace{c^0_c+c^{u\top}_i c^c_c+c^{s\top}_c x_{ict}}_{\text{household and timing}}
+\underbrace{\kappa_c\bigl(IV_{ict}-\overline{IV}_c\bigr)}_{\text{how good the items are}} ,
\label{eq:inc}
\end{equation}
with $\sigma(x)=1/(1+e^{-x})$ turning the score into a probability. Term by term:

\begin{center}
\begin{tabular}{lll}
\toprule
symbol & shape & what it is \\
\midrule
$y_{ct}$ & scalar $\in\{0,1\}$ & 1 if the household buys anything from category $c$ on this trip \\
$c^0_c$ & one per category & how often $c$ is bought at all: bread often, motor oil rarely \\
$c^u_i$ & $K$ per household & what kinds of category this household leans toward \\
$c^c_c$ & $K$ per category & where category $c$ sits among those kinds \\
$c^{u\top}_i c^c_c$ & scalar & their match: this household's appetite for this category \\
$x_{ict}$ & $4$ numbers & the recency basis of Section~\ref{sec:util}, for this category \\
$c^s_c$ & $4$ per category & how strongly recency matters for this category \\
$\kappa_c$ & one per category & how much the items' appeal feeds into buying the category \\
$IV_{ict}$ & scalar & how good the category's items look today (Section~\ref{sec:nest}) \\
$\overline{IV}_c$ & one per category & a constant subtracted from $IV$, fixed once during \\
 & & fitting (Section~\ref{sec:iv}) \\
\bottomrule
\end{tabular}
\end{center}

So the first three terms answer ``would this household buy this category on a normal
day?'' and the fourth adjusts for what the category is worth \emph{today}, given its
prices and what the store has.

\paragraph{What $x_{ict}$ actually contains, and why it is wrong.} The recency basis of
Section~\ref{sec:util} is defined per \emph{sub-commodity}. A category contains several,
so a category-level recency has to be constructed somehow. The code does this
(\code{27\_nested\_basket.py:425}):

\begin{enumerate}
\item take the item in category $c$ with the lowest \code{item\_id}, which is the lowest
\code{PRODUCT\_ID} --- an arbitrary ordering with no meaning;
\item read off that one item's sub-commodity;
\item set $\tau$ = days since household $i$ last bought \emph{anything in that
sub-commodity};
\item return $x_{ict}=\bigl[\mathbb 1\{\text{never}\},\ e^{-\tau/7},\
e^{-\tau/g_{\mathrm{sub}}},\ \log(1+\tau)/\log 100\bigr]^{\top}$.
\end{enumerate}

So it is not a category-level quantity at all. It is the recency of one arbitrarily chosen
sub-commodity, standing in for the category. Concretely:

\begin{center}
\small\begin{tabular}{lrrlr}
\toprule
category & items & subs & sub-commodity actually used & \% of items \\
\midrule
BAKED BREAD/BUNS/ROLLS & 141 & 14 & HOT DOG BUNS & \textbf{8.5}\% \\
BAG SNACKS & 153 & 12 & PRETZELS & 12.4\% \\
CHEESE & 176 & 11 & NATURAL CHEESE CHUNKS & 19.9\% \\
FRZN MEAT/MEAT DINNERS & 201 & 7 & PREMIUM ENTREES & 21.9\% \\
SOFT DRINKS & 225 & 14 & 2 LITER BOTTLE, CARBONATED & 30.7\% \\
YOGURT & 148 & 2 & YOGURT NOT MULTI-PACKS & 76.4\% \\
\bottomrule
\end{tabular}
\end{center}

For bread, the model's answer to ``is this household due to buy bread?'' is ``have they
recently bought hot dog buns?''. Across all 188 categories the median chosen
sub-commodity covers 50\% of its category's items, and in 21\% of categories it
covers under 20\%.

This is a defect, not a modelling choice, and we record it as one (Section~\ref{sec:disc}).
Its effect is bounded by what the term is worth: $c^{s\top}_c x_{ict}$ sits inside the
incidence score alongside the household and inclusive-value terms, and the state ablation
in Section~\ref{sec:ablate} bounds the whole recency mechanism --- item-level and
category-level together --- at 0.071 nats. A correct version would use the days since
the household last bought \emph{any} item in the category, which is a change to the data
build rather than the model; we have not made it, so every number in this paper uses the
arbitrary-sub-commodity version.

\paragraph{2. How many different products, given it is bought?}
\begin{equation}
k_{ct}-1\sim\mathrm{Poisson}\bigl(e^{z^b_{ict}}\bigr),\qquad
z^b_{ict}=b^0_c+b^u_i-b^p_c\,\overline{\Delta\log p}_{ct},
\label{eq:brd}
\end{equation}
Subtracting 1 and adding it back means $k\ge1$: a bought category contains at least one
product, and the Poisson only models how many \emph{extra} ones.

\begin{center}
\begin{tabular}{lll}
\toprule
symbol & shape & what it is \\
\midrule
$k_{ct}$ & integer $\ge1$ & how many different products are taken from category $c$ \\
$b^0_c$ & one per category & the typical breadth of this category --- one yoghurt or four \\
$b^u_i$ & one per household & whether this household buys broadly or narrowly, overall \\
$b^p_c$ & one per category & how much a category-wide promotion widens the selection \\
$\overline{\Delta\log p}_{ct}$ & scalar & average price deviation over the category's stocked items \\
\bottomrule
\end{tabular}
\end{center}

The minus sign on $b^p_c$ means a positive $b^p_c$ widens the basket when prices fall.

\paragraph{3. Which products?} The $k_{ct}$ products are drawn, without replacement, in
proportion to $e^{u}$ among the items the store carries --- so an item twice as appealing
is twice as likely, and an item the store does not stock cannot be chosen at all:
\begin{equation}
P(j)=\frac{e^{u_{ijt}}}{\sum_{j'\in\mathcal{C}_s(c)}e^{u_{ij't}}},\qquad j\in\mathcal{C}_s(c).
\label{eq:alloc}
\end{equation}

\paragraph{4. How many units of each chosen product?}
\begin{equation}
q_{jt}-1\sim\mathrm{Poisson}\bigl(e^{z_{ijt}}\bigr),\qquad
z_{ijt}=q^0_j-(\gamma^{q\top}_i\beta^q_j)\,\Delta\log p_{jst}+\eta^{q\top}_j x_{ijt},
\label{eq:qty}
\end{equation}
Again $q\ge1$: a chosen product is bought at least once, and the Poisson models the
extras.

\begin{center}
\begin{tabular}{lll}
\toprule
symbol & shape & what it is \\
\midrule
$q_{jt}$ & integer $\ge1$ & units of product $j$ \\
$q^0_j$ & one per item & how many people usually take: one ketchup, six yoghurts \\
$\gamma^q_i$ & $K_p$ per household & willingness to stock up when something is cheap \\
$\beta^q_j$ & $K_p$ per item & how stockpile-able this item is \\
$\eta^q_j$ & $4$ per item & how time since last purchase changes the amount \\
\bottomrule
\end{tabular}
\end{center}

Note the price coefficient here, $\gamma^{q\top}_i\beta^q_j$, is \emph{separate} from the
one in Eq.~\ref{eq:util}. That is deliberate: the data says
whether a household buys an item and how much of it respond to price at different rates
(Section~\ref{sec:elast}), and sharing one coefficient would assume away the difference.

\subsection{The term that ties them together, and what \texorpdfstring{$\kappa$}{kappa} means}
\label{sec:nest}

The four answers would be independent were it not for one term in Eq.~\ref{eq:inc}. The
\emph{inclusive value}
\begin{equation}
IV_{ict}=\log\!\!\sum_{j\in\mathcal{C}_s(c)}\!\! e^{u_{ijt}}
\label{eq:iv}
\end{equation}
collapses a whole category into one number: how appealing its items are, right now, at
this store, at today's prices. If everything in the category goes on offer, every
$u_{ijt}$ rises and so does $IV$. The coefficient $\kappa_c$ decides how much of that
reaches the decision to buy the category at all.

This is the model's answer to the retailer's question. A price cut on one item raises
$IV$ a little; $\kappa_c$ says whether that turns into new category buyers or merely
reshuffles Eq.~\ref{eq:alloc}:

\begin{center}
\begin{tabular}{ll}
\toprule
$\kappa_c=0$ & incidence ignores the items completely. A price cut moves share inside the
category \\ & and nothing else --- pure cannibalisation. \\
$\kappa_c=1$ & the reference point (see below): the category grows exactly as much as its
items \\ & gained, and no more. \\
$\kappa_c>1$ & the category grows more than its items gained --- the promotion pulls in
trips. \\
\bottomrule
\end{tabular}
\end{center}

\paragraph{Why $\kappa=1$ is the reference point.} Suppose instead of a coin flip we
modelled the total units a household takes from category $c$ as a count
$Q\sim\mathrm{Poisson}(e^{a_{ic}+\kappa_c IV_{ict}})$, split across the category's items in
proportion to $e^{u}$. A Poisson total split by fixed proportions gives independent
Poisson counts per item, with rates $e^{a_{ic}+\kappa_c IV}\cdot e^{u_{ijt}-IV}$, because
the proportion for item $j$ is exactly $e^{u_{ijt}-IV}$ by Eq.~\ref{eq:iv}. That is
\begin{equation}
q_{ijt}\sim\mathrm{Poisson}\bigl(e^{\,a_{ic}+(\kappa_c-1)IV_{ict}+u_{ijt}}\bigr),
\label{eq:kappa1}
\end{equation}
and at $\kappa_c=1$ the inclusive value disappears: each item's rate depends only on its
own $u$, so the category's total is whatever its items happen to add up to. That is the
sense in which $\kappa=1$ means ``no expansion'', and it is why we parameterise
$\kappa_c=\mathrm{softplus}(\kappa^{\mathrm{raw}}_c)$ starting at exactly $1.0$: the data
must push it off the neutral point in one direction or the other.

Eq.~\ref{eq:kappa1} is a device for reading $\kappa$, not a second model. The fitted
model is Eqs.~\ref{eq:inc}--\ref{eq:qty}, in which the inclusive value appears once, in
the incidence term.

\subsection{Letting items in a cart influence each other}
\label{sec:inter}

The remaining term in Eq.~\ref{eq:util} is $\rho_j^{\top}\bar\alpha_j(S)$, where $S$ is
the set of items in the basket. For item $j$ in a basket of $n$ items, the context is
\begin{equation}
\bar\alpha_j(S)=\frac{1}{n-1}\Bigl(\sum_{k\in S}\alpha_k-\alpha_j\Bigr),
\qquad \bar\alpha_j(S):=\mathbf 0 \text{ when } n=1.
\label{eq:ctx}
\end{equation}

\paragraph{Why this is an average over the other items.} Item $j$ is itself in the basket,
so $\sum_{k\in S}\alpha_k$ already includes $\alpha_j$; subtracting it leaves the sum over
everything else,
\begin{equation}
\sum_{k\in S}\alpha_k-\alpha_j=\sum_{k\in S,\,k\neq j}\alpha_k ,
\qquad\text{so}\qquad
\bar\alpha_j(S)=\frac{1}{n-1}\sum_{k\in S,\,k\neq j}\alpha_k .
\end{equation}
$S$ has $n$ items and one of them is $j$, so there are $n-1$ others; dividing by $n-1$
makes this their average. When $n=1$ there are no others, $n-1=0$, and the term is set to
zero.

\paragraph{Why the interaction is an average of pairwise affinities.} A dot product is
linear, so the constant comes out and $\rho_j^{\top}$ goes inside the sum:
\begin{equation}
\rho_j^{\top}\bar\alpha_j(S)=\frac{1}{n-1}\sum_{k\in S,\,k\neq j}\rho_j^{\top}\alpha_k .
\label{eq:pairwise}
\end{equation}
Each $\rho_j^{\top}\alpha_k$ is one number saying how much having $k$ in the cart helps or
hurts $j$: positive means they go together, negative means they compete. The term is the
average of those over everything else in the cart. For a cart $\{j,k,l,m\}$ it is
$\tfrac13(\rho_j^{\top}\alpha_k+\rho_j^{\top}\alpha_l+\rho_j^{\top}\alpha_m)$.

Eq.~\ref{eq:ctx} is what the code computes, because the basket sum is accumulated once and
each item's context follows from one subtraction. The data records no order within a
receipt, so $S$ is a set: the context is the whole cart, never a prefix.

\paragraph{Tied or untied.} Setting $\rho_j=\alpha_j$ makes the interaction
$\alpha_j^{\top}\alpha_k$, which is symmetric and largest when two items are \emph{alike}.
That forces co-purchase structure into $\alpha$, which is what makes the embedding recover
product categories (Section~\ref{sec:repr}), but it cannot say that two \emph{different}
items belong together. Leaving $\rho$ free allows that. The tied form is nevertheless the
default, and not only by preference: Eq.~\ref{eq:swap} requires the pairwise interaction to
be \emph{symmetric}, and an untied $\rho_j^{\top}\alpha_k$ is not, so no potential function
exists and $P(S\mid\mathcal K)$ would not be defined. Symmetry can be had without tying ---
$\tfrac12(\alpha_j^{\top}\rho_k+\rho_j^{\top}\alpha_k)$ admits a potential --- so tying is
a stronger restriction than the theory needs, and it costs the ability to represent an
arbitrary symmetric co-purchase structure: $W=AA^{\top}$ is positive semidefinite of rank at
most $K$.

\paragraph{What distribution over baskets this implies.} Write $b_j$ for every part of
Eq.~\ref{eq:util} that does not involve other basket items. For baskets of a fixed size
$n$, define a score for the whole basket,
\begin{equation}
E(S)=\sum_{j\in S}b_j+\frac{1}{n-1}\sum_{\{j,k\}\subset S}\alpha_j^{\top}\alpha_k ,
\label{eq:energy}
\end{equation}
where $\{j,k\}\subset S$ runs over each unordered pair once, and let higher-scoring
baskets be more likely, $P(S)\propto e^{E(S)}$. Now swap one item: replace $j\in S$ by
some $m\notin S$, giving $S'$. The size does not change, so the constant $\tfrac1{n-1}$ is
the same on both sides. Every item other than $j$ and $m$ contributes the same $b$ to both
and cancels; every pair not involving $j$ or $m$ likewise cancels. What is left is the
pairs $j$ had with the rest, removed, and the pairs $m$ gains with that same rest:
\begin{equation}
E(S')-E(S)=(b_m-b_j)+\alpha_m^{\top}\bar\alpha_j(S)-\alpha_j^{\top}\bar\alpha_j(S)=u_m-u_j .
\label{eq:swap}
\end{equation}
That is exactly the difference Eq.~\ref{eq:alloc} scores when choosing between $m$ and
$j$. So the model's answer to ``what goes in this slot, given the rest of the cart'' is
the conditional of $P(S)\propto e^{E(S)}$ at fixed basket size --- which is what makes the
Gibbs sampler of Section~\ref{sec:sampling} legitimate, and is checked numerically in
Section~\ref{sec:verify}.

\subsection{The probability of a whole basket}
\label{sec:joint}

Multiplying the four answers together gives what the model says about a trip. Write
\begin{equation}
\mathcal{B}=\bigl(\;\underbrace{y_1..y_C}_{\text{which categories}},\;
\underbrace{(k_c)}_{\text{how many products each}},\;
\underbrace{(J_c)}_{\text{which products}},\;
\underbrace{(q_j)}_{\text{how many units each}}\;\bigr)
\end{equation}
for a basket, where $J_c$ is the set of $k_c$ products taken from category $c$. Prices,
the store, the day and the household are given, never generated. Then
\begin{equation}
P(\mathcal{B})=
\prod_{c=1}^{C}\underbrace{\sigma(\eta_{ict})^{y_c}\bigl(1-\sigma(\eta_{ict})\bigr)^{1-y_c}}_{\text{Eq.~\ref{eq:inc}: bought or not}}
\;\cdot\!\!\prod_{c:\,y_c=1}\!\!\Bigl[\;
\underbrace{P(k_c)}_{\text{Eq.~\ref{eq:brd}}}\;
\underbrace{P(J_c\mid k_c)}_{\text{Eq.~\ref{eq:alloc}}}\;
\prod_{j\in J_c}\underbrace{P(q_j)}_{\text{Eq.~\ref{eq:qty}}}\Bigr].
\label{eq:joint}
\end{equation}
Every category contributes a coin-flip factor whether bought or not; a bought category
contributes three more. Reading Eq.~\ref{eq:joint} left to right is exactly the sampler of
Section~\ref{sec:sampling}, stage by stage.

\subsection{What is fitted, and the three ways it differs from Eq.~\ref{eq:joint}}
\label{sec:fitgap}

Training does not maximise $\log P(\mathcal{B})$. It minimises a weighted sum of four
losses (Section~\ref{sec:obj}), and three differences are worth naming because they change
what the reported numbers mean.

\emph{(i) The choice term is conditional on the rest of the cart.} Eq.~\ref{eq:alloc}
draws $k_c$ products with no reference to the other categories. Training instead scores
each purchased row against the whole rest of the basket, through $\bar\alpha_j(S)$, and so
maximises $\sum_j\log P(j\mid\mathcal{B}\setminus j)$. Eq.~\ref{eq:swap} shows those
conditionals belong to a genuine distribution, which is what licenses the sampler; but
they are not the factor $P(J_c\mid k_c)$ in Eq.~\ref{eq:joint}.

\emph{(ii) The two choice sets once differed, and no longer do.} Earlier versions trained
the item head against negatives drawn from the \emph{whole catalogue} while generating
within a category, so the term was fitted as a catalogue-wide choice and used as a
within-category one. Eq.~\ref{eq:alloc} now normalises over the category's stocked items in
both places (Section~\ref{sec:negsamp}), so this particular gap is closed. It is recorded
because the co-occurrence results in Section~\ref{sec:gen} were once attributed to it.

\emph{(iii) The terms are separately averaged and weighted.} Each loss is a mean over its
own sample --- purchase rows, sampled categories, bought categories --- then multiplied by
$w_q,w_c,w_b$. Eq.~\ref{eq:joint} implies particular relative weights; the objective does
not reproduce them.

\paragraph{Consequence.} The model has a well-defined distribution over baskets,
Eq.~\ref{eq:joint}, and Section~\ref{sec:gen} tests it by sampling from it and comparing
against real baskets. It has no tractable likelihood for a whole basket: no number in this
paper is $\log P(\mathcal{B})$ on held-out data. Every likelihood reported is the
conditional in (i).

%======================================================================
\section{Inference}
\label{sec:inf}

We fit by stochastic gradient descent on a weighted sum of the four terms. This section
gives the objective, the three sampling schemes it depends on, and the computational
cost. Code: \code{scripts/27\_nested\_basket.py}.

\subsection{Objective}
\label{sec:obj}

\begin{equation}
\mathcal L=\mathcal L_{\text{item}}+w_q\mathcal L_{\text{qty}}
+w_c\mathcal L_{\text{inc}}+w_b\mathcal L_{\text{brd}}
-\tau_{\text{repr}}\lVert\Theta_{\text{repr}}\rVert^2
-\tau_{\text{inc}}\lVert\Theta_{\text{inc}}\rVert^2
-\tau_{\text{price}}\lVert\Theta_{\text{price}}\rVert^2 .
\label{eq:obj}
\end{equation}
\textbf{Three} penalty blocks, not two. $\Theta_{\text{repr}}$ covers
$\alpha,\theta,\eta,\mu,\delta,\zeta,\xi,w^{\mathrm{loy}},w^{\mathrm{frq}}$;
$\Theta_{\text{inc}}$ covers $c^u,c^c,c^s$; $\Theta_{\text{price}}$ covers only
$\gamma,\beta,\gamma^q,\beta^q$, at $10^{-4}/1500$ against $10^{-2}/1500$ for the
representation block. Shrinking the price block at the representation block's rate biases
the elasticity toward zero, and the elasticity is the estimand.

$\Theta_{\text{inc}}$ has to be separate, and this is not a tuning preference. Fitted inside
$\Theta_{\text{repr}}$ under a single $\tau$, the incidence embeddings collapse to near rank
one: 144{,}256 parameters delivering roughly one number per household times one per
category, which cannot express ``buys nappies, never cat food''. Averaged over all
dimensions the data gradient beats the penalty 5.8--17$\times$, but that average is carried
by the one direction that is learning; on the rest the data gradient is near zero and the
penalty is the only force. Separating the block raises the entropy effective rank of $c^u$
from 5.2 to 31.8 and of $c^c$ from 4.1 to 23.7, and improves incidence NLL
$0.1249\to0.1199$ against a $0.1132$ oracle floor.

\paragraph{A discrepancy worth stating.} Eq.~\ref{eq:obj} is not quite what the code
maximises. The $C/|\mathcal S_v|$ rescaling that makes $\mathcal L_{\text{inc}}$ an unbiased
estimate of the full 188-category likelihood means that at $w_c=1$ the incidence term
carries $M\cdot C\approx 36{,}096$ summands against the item term's ${\approx}1{,}476$; the
implementation instead takes a plain mean of each component over its own term count, which
weights incidence about $24.5\times$ lower. Fitting Eq.~\ref{eq:obj} literally
($w_c=24.46$) buys $0.0035$ nats of incidence likelihood and costs \textbf{63\% of the price
coefficient}, $+0.695\to+0.260$ --- roughly 42 refit standard deviations --- because
$\gamma,\beta$ reach incidence only through $b\to IV$, while the item head carries 79\% of
the own-price elasticity. Neither weighting is \emph{a priori} correct, since any $w\neq1$
forfeits the pseudo-likelihood consistency argument; every number in this paper comes from
the mean-over-terms convention.

Optimiser Adam, learning rate 0.005 annealed by cosine to $0.05\times$ over 6{,}000
iterations. A batch is 192 baskets expanded to all their purchase rows. Validation runs
every 1{,}000 iterations and the checkpoint is kept on the sum of all four held-out terms;
because the item conditional is exact, that criterion needs no surrogate and no
calibration.

\subsection{There is no negative sampling}
\label{sec:negsamp}

Earlier versions of this model sampled $K$ negatives from
$p(j)\propto\max(\mathrm{count}_{\text{train}}(j),1)^{0.75}$ over the whole catalogue and
scored a softmax over $\{$chosen$\}\cup\{$negatives$\}$. That is no longer done anywhere,
and the reason is structural rather than computational.

The conditional the model actually defines normalises over $A_{-j}$, the purchased item's
own category minus whatever else the basket holds from it --- a median of 15 products and at
most 225. That sum is affordable \emph{exactly}. So the item head evaluates the true
conditional, and with it go the proposal distribution, the $\log q$ correction, and any
metric whose value depends on how decoys were drawn.

This matters for reading the literature as much as for the fit. A catalogue-wide decoy is
almost always separable by $\lambda_j+\theta_i^{\top}\alpha_j$ alone, so a sampled loss asks
``is this item plausible at all'' rather than ``does this item fit \emph{this} basket'', and
its optimal scorer is $\log p-\log q$ rather than $\log p$. Item log-likelihoods reported
against sampled decoys --- including those in earlier versions of this paper --- are not
comparable to the numbers in Section~\ref{sec:emp}; only distance from chance carries
across.

\subsection{The inclusive value is estimated, not summed}
\label{sec:iv}

$IV_{ict}$ is a log-sum-exp over a category's stocked items. Categories are padded to the
largest (225 items) although the median has 15, so a dense evaluation spends
most of its work on padding. Write $n_c$ for the number of items of category $c$ the trip's store carries. We sample
$m=32$ of them uniformly with replacement and rescale:
\begin{equation}
\widehat{IV}=\log\Bigl(\sum_{s=1}^{m}e^{u_{j_s}}\Bigr)+\log\frac{n_c}{m}.
\end{equation}
The \emph{sum} inside is unbiased, since
$\mathbb E[\frac{n_c}{m}\sum_s e^{u_{j_s}}]=\frac{n_c}{m}\cdot m\cdot\frac{1}{n_c}\sum_j e^{u_j}
=\sum_j e^{u_j}$. Its \emph{logarithm} is not: by Jensen
$\mathbb E[\log X]\le\log\mathbb E[X]$, so $\widehat{IV}$ is biased slightly low. The bias
is absorbed into the category intercept $c^0_c$ during training, and generation uses the
same estimator, so the two remain consistent.

\paragraph{A frozen reference.} $\overline{IV}_c$ in Eq.~\ref{eq:inc} is fixed once, at
iteration 500, to the mean $IV$ over a \emph{uniform} category sample. Any fixed
constant is absorbed by $c^0_c$, so freezing does not change the fit; it guarantees that
training and generation subtract the same thing. Taking the reference from the training
batches instead uses whatever category mix those batches had, and if that mix differs
from the population the generated baskets come out systematically wrong.

\subsection{Category sampling for the incidence head}

Incidence is a Bernoulli over 188 categories per trip, of which only 6.12 are
bought. We sample 16 categories per trip \textbf{uniformly}. Case-control sampling
(a few bought, a few not) is tempting at a 3.25\% base rate, and biases the fit two
ways: the intercept, which the standard $\log(\pi_1/\pi_0)$ offset corrects, and the term
$\kappa_c(IV-\overline{IV}_c)$, whose mean differs between sample and population because
bought categories have higher inclusive values --- which the offset does not touch.
Uniform sampling removes both at source and needs no correction.

\subsection{Computational cost}

\begin{center}
\begin{tabular}{lrrr}
\toprule
run & iterations & wall time & ms/iter \\
\midrule
default & 12,000 & 10.8 min & 53.8 \\
\bottomrule
\end{tabular}
\end{center}

Two lookups dominate and both are sparse and high-cardinality: household state and
store price. Materialising (household $\times$ day $\times$ sub-commodity) is about 32
million rows and a per-sample dictionary is about 35 million hits per epoch. Both are
instead stored once as a globally sorted key array with
$\mathrm{key}=\mathrm{group}\cdot\mathrm{stride}+\mathrm{index}$, so a whole batch
resolves in one \code{np.searchsorted}; the stride exceeds any index so ordering never
crosses a group boundary. That search is still roughly 40\% of a step, and is avoidable
by deduplicating repeated candidates within a batch, which we have not done.

\paragraph{Scalability.} Cost is independent of catalogue size (negative sampling fixes
the candidate count), of the number of households (only the embedding table grows), and
of the number of baskets per step. Categories are capped at 16 sampled per trip and
items per category at 32 by the sampled inclusive value. Nothing is quadratic in
items or households.

%======================================================================
\section{Sampling}
\label{sec:sampling}

Generation is a harder requirement than prediction and is where most of this paper's
findings came from. We state precisely what the model gives us and what it does not.

\paragraph{What we have.} Conditionals: $P(\text{item}\mid\text{rest of basket})$ from the
item head, $P(\text{category bought})$ from the incidence head, and count distributions
from the quantity and breadth heads. \textbf{What we do not have} is a normalised joint
$P(S)$: computing it needs a partition function over all size-$n$ baskets, which is
intractable. Eq.~\ref{eq:swap} shows the item head's conditionals are those of
$P(S)\propto e^{E(S)}$ at fixed $n$, so we can \emph{sample} that joint without ever
evaluating it.

\begin{algorithm}[h]
\SetAlgoLined
\KwIn{household $i$, day $t$, week $w$, store $s$}
\KwOut{a basket of items with unit counts}
\textbf{Stage 1 --- categories.} For each $c$, compute $IV_{ict}$ and draw
$\mathbb 1\{c\in B\}\sim\mathrm{Bernoulli}(\sigma(\mathrm{logit}_{ict}))$\;
\If{the category interaction is on}{
  Gibbs sweep the $C$-dimensional indicator: for each $c$, recompute
  $\mathrm{logit}_{ict}$ given the current $B\setminus c$ and redraw\;
}
\textbf{Stage 2 --- how many.} For each $c\in B$ draw
$k_c=1+\mathrm{Poisson}(e^{z^b_{ic}})$, capped at the number the store stocks\;
\textbf{Stage 3 --- which.} For each $c\in B$, draw $k_c$ distinct items from
$\mathrm{softmax}_j u_{ijt}$ over that category's stocked items, with context zero (no
basket exists yet)\;
\textbf{Stage 4 --- refine.} Gibbs sweeps: for each filled slot, remove its item,
recompute $\bar\alpha$ from the current draft, and redraw from that slot's own category\;
\textbf{Stage 5 --- units.} For each chosen item draw $1+\mathrm{Poisson}(e^{z_{ijt}})$\;
\caption{Basket generation}
\end{algorithm}

Stage 4 holds basket size and category composition fixed, which is the move
Eq.~\ref{eq:swap} is defined over. Stage 1's optional sweep is the same argument one
level up: the category indicator is a $C$-dimensional binary field whose single-site
conditionals the incidence head supplies.

\paragraph{A sequential alternative.} If the model is trained with a random-prefix
context --- one permutation per basket, each row seeing only the items before it, which
is SHOPPER's unordered likelihood (its Eq.~6) --- then Stages 3--4 collapse into one
pass that places items one at a time with $\bar\alpha$ the running mean of what is
already placed. Training and generation then use identical conditionals and no Gibbs is
needed. This is implemented behind a flag and is not evaluated in this paper.

\paragraph{What generation does not carry.} With the context zeroed at Stage 3 and no
Stage 4, generated baskets reproduce marginals only. This is not a detail: it is what
Section~\ref{sec:gen} measures.

%======================================================================
\section{From measurement to model}
\label{sec:mapping}

\begin{center}
\begin{tabular}{p{0.42\textwidth}p{0.5\textwidth}}
\toprule
what Section~\ref{sec:eda} showed & what the model does \\
\midrule
25\% of the units response is existing buyers taking more &
quantity term with its own price coefficient \\
73\% of the units response is entering a category at all &
the nest, with $\kappa_c$ estimated rather than assumed \\
1.284 distinct items per purchased category; 22.3\% of rows buy more than one unit &
breadth term; units kept as counts, never binarised \\
Taste is 1.66$\times$ more self-similar within household than across &
per-household $\theta_i$ \\
Per-household price slopes have split-half reliability $+0.236$ &
price sensitivity low-rank $\gamma_i^{\top}\beta_j$, not free per pair \\
Median store carries only 63\% of the catalogue &
availability mask: unstocked items leave the choice set \\
15.8\% of store-item-weeks differ from the chain price &
store-level price deviation where observed \\
Week effects carry a confound between price and demand &
$\mu_j^{\top}\delta_w$ seasonality control \\
Promotion weeks are $3.4\times$ steeper than ordinary variation &
\emph{not addressed}; Section~\ref{sec:disc} \\
\bottomrule
\end{tabular}
\end{center}

%======================================================================
\section{Empirical study}
\label{sec:emp}

Every number in this section comes from one fitted model. That sentence is doing work: an
earlier version of this section mixed results from two fits without naming either, and two
of its conclusions did not survive being recomputed on a single one.

\subsection{Protocol, and what is comparable}
\label{sec:protocol}

All numbers are on held-out test weeks. The item score is the exact conditional --- a
softmax over the purchased item's category minus whatever else the basket holds from it ---
so it depends on no sampling scheme. The choice set averages \textbf{46.9} products, making
chance $-\log 46.9=-3.848$ and $2.1\%$ top-1.

Two sources of uncertainty are separated because they cover different things. \textbf{Refit
variability} comes from seven independent fits differing only in seed: item sd
\textbf{0.0040}, top-1 sd 0.0013, price coefficient sd 0.0104, $\kappa$ sd \textbf{0.0011}.
A gap between two independently fitted models must exceed \textbf{0.0110} nats to clear it
at 95\%. An eighth fit at a seed outside the seven landed at item $-2.5024$, price $+0.7085$,
$\kappa$ $0.8701$ --- inside the distribution on all three, which is the only out-of-sample
check these standard deviations have. \textbf{Evaluation sampling} is a household block
bootstrap: 400 draws over 1{,}664 households and 6{,}000 baskets, paired so the interval is
on the difference. Neither covers resampling the training households, which would need a
refit per replicate.

\subsection{Do the equations match the code?}
\label{sec:verify}

Thirty checks compare each printed formula against the fitted model --- by autograd where
this paper states a derivative, by direct evaluation where it states an identity, and by
rebuilding the quantity from the raw training split where it states a definition
(\code{eval/33\_verify\_equations.py}). All thirty pass, on every model fitted
for this paper.

They did not always. An earlier version of this paper claimed fifteen checks passing while
the suite could not run at all: its label argument was never read, so it silently loaded a
superseded checkpoint. Behind that, its quantity check still indexed candidate slot~0 after
the loss was fixed to index the purchased item, so it had been \emph{failing}; and its
term-by-term check for Eq.~\ref{eq:util} omitted two terms and passed anyway, because the
pair it tested happened to have no purchase history. Three divergences between this paper
and the code surfaced once the suite ran --- the persistence terms of Eq.~\ref{eq:persist},
the habit loading in the incidence predictor, and the third penalty block of
Eq.~\ref{eq:obj} --- and all three are now stated here because the code has always had them.

\subsection{Ablations}
\label{sec:ablate}

Eight fits, differing only in which term is removed, all 6{,}000 iterations at one seed with
identical data and schedule.

\begin{center}
\begin{tabular}{lrrl}
\toprule
removing & item log-lik & cost & clears refit noise? \\
\midrule
--- (full model) & $\mathbf{-2.5118}$ & --- & --- \\
quantity head & $-2.5113$ & $-0.0005$ & \textbf{no}, and correctly so \\
nest & $-2.5334$ & $+0.0215$ & yes, by $2\times$ \\
household state & $-2.5373$ & $+0.0255$ & yes \\
prices scrambled & $-2.5706$ & $+0.0587$ & yes \\
basket interaction & $-2.5818$ & $+0.0700$ & yes \\
persistence & $-2.5833$ & $+0.0714$ & yes \\
store & $-2.8240$ & $+0.3122$ & yes \\
\bottomrule
\end{tabular}
\end{center}

With bootstrap intervals on the paired gaps: no persistence $+0.0728$ $[+0.0648,+0.0816]$;
the shared incidence penalty $+0.0715$ $[+0.0635,+0.0804]$; prices scrambled $+0.0562$
$[+0.0489,+0.0626]$. Every interval excludes zero.

Store dominates at $0.312$, nearly four times the next term. Removing the quantity head
costs $-0.0005$ --- inside refit noise and marginally the wrong way, which is what a null
looks like, and the only evidence here that the procedure can return one. $\kappa$ is stable
at $0.865$--$0.872$ across every variant including the placebo, and sits below 1 throughout.

\paragraph{Two terms that look independent are substitutes.} Before the persistence terms of
Eq.~\ref{eq:persist} were added, the household state was worth $+0.049$ and the basket
interaction $+0.058$. State has now halved to $+0.026$ while the interaction grew to
$+0.070$. $\omega_{ij}$ and $f_{ij}$ are per-(household, item) purchase history, which
overlaps with what $x_{ijt}$ carried: persistence did not simply add signal, it absorbed
about half of the state term's job.

\paragraph{Removing the nest is the only ablation that lowers the price coefficient}
($+0.695\to+0.596$). With no $\kappa\cdot IV$ channel, price cannot reach category entry at
all, and the fit shrinks the coefficient rather than redistributing it. Every other ablation
nudges it up, toward $+0.71$--$0.75$.

\subsection{Against learned baselines}
\label{sec:strong}

HPF and B-Emb are refitted and scored on identical choice sets. They bracket the model
deliberately: HPF has the household half and no basket interaction, B-Emb the basket half
and no household.

\begin{center}
\begin{tabular}{lrr}
\toprule
model & log-lik & top-1 \\
\midrule
HPF, Poisson factorisation & $-3.1866$ & 0.165 \\
B-Emb, skip-gram embeddings & $-3.0939$ & 0.187 \\
full model, no persistence & $-2.5787$ & 0.322 \\
full model & $\mathbf{-2.5080}$ & \textbf{0.341} \\
\bottomrule
\end{tabular}
\end{center}

The gap over the stronger baseline is $+0.5860$ nats, 95\% CI $[+0.5628,+0.6094]$ by paired
household bootstrap.

\paragraph{SHOPPER, on both metrics.} SHOPPER ships as CUDA C++ and could not be run on the
available hardware, so it is reimplemented in PyTorch from the paper and its source: untied
$\rho$, sequential prefix context, catalogue-wide softmax. Its choice set is the whole
catalogue while ours is a category, so scoring only on ours would flatter us; both are
therefore scored on both, exactly. Within category SHOPPER reaches $-2.8262$ against our
$\mathbf{-2.5058}$; on the full catalogue $-7.0501$ against $\mathbf{-6.7402}$. The margin of
$+0.310$ nats survives on SHOPPER's own normalisation, which is the comparison that matters.

\paragraph{This is a floor, not a fair fight.} Same metric and choice sets, but one tuned
scalar per baseline against 881{,}490 parameters and an architecture selected over many
runs. SHOPPER is also Bayesian and maximises an ELBO while this fits the same model by MAP,
so the comparison is between models, not inference procedures, and the SHOPPER column
carries no uncertainty of its own.

\subsection{Is the price response identified?}
\label{sec:causal}

Refitting on a price panel scrambled before fitting, store deviations included:

\begin{center}
\begin{tabular}{lrrr}
\toprule
model & price coef. & allocation elasticity & $\kappa$ \\
\midrule
fitted & $+0.695$ & $-0.7171$ $[-0.7646,-0.6713]$ & 0.8723 \\
prices scrambled & $\mathbf{+0.000}$ & $\mathbf{+0.0000}$ $[-0.0029,+0.0000]$ & 0.8720 \\
\bottomrule
\end{tabular}
\end{center}

The placebo retains 0.0\% of the coefficient on both margins --- the fitted value is
$5\times10^{-39}$ --- while $\kappa$ moves by $0.0003$ and the incidence likelihood is
untouched. Prices scramble; category structure does not.

\paragraph{What this does and does not show.} Permuting a product's price series across days
destroys the price--demand association \emph{whether it is causal or confounded}. It rules
out a spurious constant. It cannot distinguish ``price moves demand'' from ``the retailer
times promotions to demand'', which is the actual endogeneity concern for within-product
price variation. There is no instrument here and no discontinuity, which is why this
subsection is no longer titled ``is the price response causal?'' --- the design does not
answer that question.

\subsection{Where a price cut goes}
\label{sec:decomp}

Own-price elasticity $-1.2070$ over 12{,}724 (household, item) pairs: allocation $-0.9562$
(79\%), incidence $-0.1393$ (12\%), quantity $-0.1115$ (9\%). Mostly a price cut moves
share. The quantity margin is unreachable for a binary-purchase model and is the concrete
return on modelling units at all. A fourth margin exists through breadth, via $-c^p_c/N_c$
in the breadth rate, and has not been measured. Median $\kappa$ is $0.872$ and \emph{no}
category exceeds 1: a price cut reaches category entry at less than full strength
everywhere.

\subsection{Generation}
\label{sec:gen}

Generated on 24{,}768 held-out trips with the same household, day, week and store.

\begin{center}
\begin{tabular}{lrrr}
\toprule
 & real & generated & TVD \\
\midrule
items per basket & 8.18 & 7.47 & 0.282 \\
categories per basket & 6.38 & 5.60 & 0.292 \\
units per basket & 11.06 & 10.23 & 0.272 \\
distinct items per category & 1.28 & 1.33 & 0.066 \\
\bottomrule
\end{tabular}
\end{center}

\textbf{The generator makes baskets too small} --- 9\% short on items and 12\% on categories.
On the superseded fit this ran the other way, so it is a defect the current model introduced
rather than an inherited one, and it is unexplained. Item marginals still correlate $+0.872$
across the catalogue with 0.4\% of real products never generated, and the price response
re-estimated from generated data recovers 74--76\% of the real elasticity. That last check is
circular --- it re-estimates a quantity the model was fitted to, so it verifies that the
sampler propagates the coefficient, not that the coefficient is right.

Co-occurrence lift correlates with real at Spearman $+0.099$ after four Gibbs sweeps, with
mean lift 3.23 real against 1.82 generated: the right ordering, weakly, and far too little
concentration. Within-basket pairs sharing a held-out sub-commodity label run $0.0482$
generated against $0.0646$ real, so fine substitution is 25\% under-produced; at department
level it is slightly \emph{over}-produced, so the failure is specific to fine distinctions.

\paragraph{A joint test.} Every likelihood above is a conditional, and comparing summary
statistics one at a time cannot detect a joint mismatch. A classifier two-sample test can: a
discriminator separates real from generated baskets at \textbf{AUC 0.7725} on eight
basket-level features. Its two most informative features are category count and basket size
--- the two marginals that are off --- so the joint failure is not yet separable from the
marginal one.

\paragraph{Novelty, and a conclusion that reversed.} Generated baskets place a
never-before-bought item 19.2\% of the time against a real 14.2\%. An earlier version of this
paper put that figure at 46\% and concluded, from a temperature sweep that barely moved it,
that the item conditional's sharpness was not the cause. Both halves were artefacts of a
measurement that attributed each generated basket to the wrong household. Measured
correctly, sharpening the generation softmax to $T=0.70$ takes novelty \emph{past} the real
rate, to 11.3\%. Sharpness is a live lever; what disqualifies it as a fix is that scaling the
item utility also scales the inclusive value, so basket size inflates from 7.55 to 12.39
items on the way.

\paragraph{Rollouts drift.} The excess compounds: cumulative distinct new items run
$1.156\times$ real after one trip and $1.261\times$ after twelve. That is a lower bound,
since generation reads the recency state from the real history, so a generated novel item
does not make the next basket more novel. Single-step counterfactuals are supported;
multi-step rollouts are usable over short horizons at most.

\paragraph{Two negative results.} Contrastive divergence with a persistent pool, which
trains directly against the model's own samples, moved the novelty rate by nothing at
either weight tested, while costing 0.03--0.04 nats. So did an explicit habit feature over household and
category, which is nonetheless still in the fitted model, at $+0.006$ nats. Both were
re-measured after the attribution bug was corrected, and both nulls stand.

\subsection{What the representations learned}
\label{sec:repr}

\paragraph{Products.} $\alpha$ recovers sub-commodity structure at \textbf{27.8$\times$}
chance on 10-nearest-neighbour purity --- $0.1126$ $[0.1041,0.1260]$ against a
permutation null at $0.0039$, a random embedding of identical shape at $0.0029$, and chance
at $0.0040$. Both nulls sit at chance, so the metric is not high for structural reasons. The
lift \emph{grows} as labels get finer: $1.3\times$ at department, $25.6\times$ at category,
$73.7\times$ at sub-commodity. But all-pairs AUC is at or below chance for the two coarse
levels ($0.496$, $0.463$) and the silhouette is negative throughout. Restricting to pairs
co-bought at least eight times, cosine separates same-sub-commodity pairs at AUC $0.928$
with mean cosine $+0.619$ against $+0.138$. $\alpha$ is organised where the data is dense
and unstructured where it is sparse.

\paragraph{Households --- a negative result that survives a non-linear probe.} Linear
\emph{and} gradient-boosted probes on $\theta$, $c^u$ and $\gamma$ fail to beat the majority
class on income, household size, age, children or home ownership; the best boosted probe
comes in $0.037$ \emph{below} it, and on income the shuffled-label control outscores the
real labels. The household parameters encode purchase behaviour, not who the household is.
At 655 households this cannot rule out a relationship a larger sample would reveal.

\paragraph{Price sensitivity.} $g_{ij}=\gamma_i^{\top}\beta_j$ has 46.8\% of its variance
between products and 17.6\% between households: price sensitivity is mostly a property of
the product, which is worth stating plainly because the household-specific slope is the
model's main claim to heterogeneity. The category ordering was never supplied and is
economically sensible --- least elastic \textsc{cigarettes} $+0.104$ and \textsc{baby foods}
$+0.127$, most elastic \textsc{bath tissues} $+2.002$ and \textsc{cold cereal} $+1.816$. The
same ten categories appear in the same two groups on the superseded fit, so this is the most
robust evidence here that $\gamma^{\top}\beta$ estimates something real. Household price
sensitivity is unrelated to income (Spearman $+0.005$, $p=0.89$), and 24.0\% of pairs carry
the economically wrong sign, which nothing in the specification forbids.

\paragraph{Coupon targeting --- a screen, not a score.} Ranking households by predicted
$g_{ij}$ does \emph{not} work: Spearman against the real held-out slope is $-0.122$, and the
quintile ordering breaks badly, with the second quintile showing a real slope of $+0.141$.
An earlier version of this paper reported a $22\times$ difference between the top and bottom
quintile pairs; on a single model the bottom two average $+0.061$, the wrong sign, and the
ratio is meaningless. What survives is the sign: pairs with $g>0$ show a real slope of
$-0.187$ against $-0.004$ for the 21\% with $g\leq0$. Used as a binary screen the
coefficient is informative; used as a graded score it is not.

\section{Related work}
\label{sec:related}

SHOPPER \cite{shopper} models a basket as a sequence of choices with an asymmetric
interaction $\rho_c^{\top}\bar\alpha$, per-customer price sensitivity as a low-rank
bilinear form, and week effects for seasonality. We inherit all three. We differ in
adding an explicit incidence stage with an estimated nesting coefficient, in modelling
unit counts rather than binary purchase, and in conditioning on the store; and we replace
its variational inference with maximum \emph{a posteriori} estimation by SGD, which loses
uncertainty quantification and gains an order of magnitude in fitting time.
Its asymmetric interaction was tried and did not improve generation on this dataset, and
its group-biased negative sampling is moot here, since the item conditional is now summed
exactly (Section~\ref{sec:negsamp}).

Section~\ref{sec:strong} compares against the two latent-factor baselines SHOPPER
uses --- hierarchical Poisson factorization \cite{hpf} and exponential family embeddings
\cite{bemb} --- and adopts its evaluation on high price-variation purchases.

The nested factorization of \cite{nf} supplies the nesting coefficient's interpretation
and the argument that week effects suffice to identify the price effect. Its
within-category softmax requires at most one item per category per trip, which
68.2\% of baskets here violate.

%======================================================================
\section{Discussion}
\label{sec:disc}

\paragraph{What the model does well.} It separates three margins that are economically
distinct and estimates their coupling; its price response survives a structural placebo;
and synthetic baskets carry 87\% of the real price elasticity, which is the
property that matters if the generator is to be used for policy learning.

\paragraph{What it does not do.}
\begin{enumerate}
\item \textbf{Per-pair co-occurrence is only partly learned.} At best the per-pair rank
correlation is $+0.091$ --- real ($z=+14.1$) but far from strong. Generated baskets have
roughly the right amount of coherence and only weakly the right pattern.
\item \textbf{Promotion weeks are steeper than a single price coefficient allows.}
Section~\ref{sec:eda} measures an implied $-2.83$ in the cut week against $-0.792$ on
ordinary variation. Display and feature are not observed, so one coefficient averages two
regimes.
\item \textbf{Basket-size distributions are wrong in shape} even where means agree.
\item \textbf{There is no held-out joint likelihood.} The model has a well-defined
generative distribution over baskets (Eq.~\ref{eq:joint}) but no tractable likelihood for
one: the item stage is fitted as a conditional, its softmax is normalised over a
different set than the generative one, and the four terms are weighted rather than
multiplied. Every likelihood in this paper is the \emph{conditional} of
Section~\ref{sec:joint}, never $\log P(\mathcal{B})$. Generation, evaluated in
Section~\ref{sec:gen}, is the only test of the joint.

\item \textbf{The generator makes baskets too small}, 9\% short on items and 12\% on
categories (Section~\ref{sec:gen}). The earlier fit erred in the opposite direction, so this
is a regression the current model introduced, and we have not explained it.
\item \textbf{$\kappa$ is weakly identified}, having moved with the incidence sampler
rather than the data before settling.
\item \textbf{Partial baskets need a correction at serving time.} The interaction term is
trained on a context averaged over a whole basket, and a one-item cart produces a vector
$2.31\times$ that norm, so without rescaling a cart with one item scores worse than an empty
one. Measured on an earlier fit and not re-measured here.

\item \textbf{No uncertainty quantification.} We report point estimates; the seed spread
bounds run-to-run noise but is not a posterior.
\end{enumerate}

\paragraph{A methodological lesson.} The most useful finding in this paper was not a
model. Eleven changes to the interaction term --- including the one a well-supported
diagnosis predicted would work --- moved the target quantity by nothing, and a change to
the negative sampling distribution moved it by fourteen standard errors. The
distinguishing evidence was a dose-response curve, and none of it would have been visible
without a model-free distributional evaluation of the generator. Where a model is meant to
generate data, it should be judged on the data it generates.

\begin{thebibliography}{9}
\bibitem{shopper}
F.~J.~R. Ruiz, S.~Athey and D.~M. Blei.
\emph{SHOPPER: A Probabilistic Model of Consumer Choice with Substitutes and Complements}.
Annals of Applied Statistics, 2020. arXiv:1711.03560.
\bibitem{nf}
S.~Donnelly, F.~J.~R. Ruiz, D.~M. Blei and S.~Athey.
\emph{Counterfactual Inference for Consumer Choice Across Many Product Categories}.
Quantitative Marketing and Economics, 2021.
\bibitem{hpf}
P.~Gopalan, J.~M. Hofman and D.~M. Blei. \emph{Scalable Recommendation with Hierarchical
Poisson Factorization}. UAI, 2015.
\bibitem{bemb}
M.~Rudolph, F.~J.~R. Ruiz, S.~Mandt and D.~M. Blei. \emph{Exponential Family Embeddings}.
NeurIPS, 2016.
\bibitem{titsias}
M.~K. Titsias. \emph{One-vs-Each Approximation to Softmax for Large Scale Classification}.
NeurIPS, 2016.
\end{thebibliography}
\end{document}
